% !PRINTED 310
\weberpage{310}
%% [page furniture, not body text -- commented out 2026-08-10] Page furniture (running
%% head, folio, signature) is excluded from this edition by standing policy.
%% \noindent Zehnter Abschnitt. \hfill \S.\,81.

%<<CONT
und da das System der $z_i'$ dieselbe Werthreihe durchläuft, wie das System der $z_i$, so ist diese Summe gleich $\Psi_{\xi_1',\xi_2',\xi_3'}$, und wir erhalten:

\gesperrt{2. Die Anwendung der Permutation $\sigma$ auf die $X$ in den Functionen $\Psi$ ist gleichbedeutend mit der Anwendung von $\varrho^{-1}$ auf die Indices $\xi_1, \xi_2, \xi_3$.}

\begin{center}
\S.\,81.\\
\gesperrt{Die Tripelsysteme der Resolventen.}
\end{center}

Wir bezeichnen jetzt die Resolvente $\Psi_{\xi_1,\xi_2,\xi_3}$ kürzer durch $\Psi^\xi$ und bemerken, dass nach der Formel (2), \S.\,80 ausser den Quadraten dieser Functionen auch noch gewisse Producte die Permutationen der Gruppe $Q$ gestatten. Zu diesen gehört zunächst das Product aller Grössen $\Psi$, das wir mit $A$ bezeichnen wollen:
\weq{1}{A = \overset{\xi}{\prod} \Psi_\xi,}
sodann aber die Producte von dreien $\Psi_\xi \Psi_\eta \Psi_\zeta$, und folglich auch ihre reciproken Werthe
\weq{2}{v = \frac{1}{\Psi_\xi \Psi_\eta \Psi_\zeta},}
wenn die Indices den Bedingungen genügen:
\weq{3}{\begin{aligned}
\xi_1 + \eta_1 + \zeta_1 &\equiv 0\\
\xi_2 + \eta_2 + \zeta_2 &\equiv 0 \quad (\mathrm{mod}\ 2).\\
\xi_3 + \eta_3 + \zeta_3 &\equiv 0
\end{aligned}}
Ein solches System von drei Functionen $\Psi$ nennen wir ein \gesperrt{Tripel}. Ebenso wollen wir drei Indexsysteme $(\xi), (\eta), (\zeta)$, die den Bedingungen (3) genügen, ein Tripel nennen.

Es giebt im Ganzen sieben und nicht mehr solcher Tripel; denn unter den drei Systemen $(\xi), (\eta), (\zeta)$ können nicht zwei einander gleich sein, und durch zwei ist das dritte eindeutig bestimmt. Man kann also für $(\xi), (\eta)$ jedes Paar aus den sieben möglichen Systemen $(\xi)$ nehmen, erhält aber dann jedes Tripel sechsmal. Die Gesammtanzahl ist also $7 . 6 : 6 = 7$.

Um die einzelnen Tripelsysteme zu charakterisiren, führen wir drei nach dem Modul 2 zu nehmende Zahlen $\alpha_1, \alpha_2, \alpha_3$ ein, die nicht alle drei verschwinden, und bemerken, dass die Congruenz
\weq{4}{\alpha_1 \xi_1 + \alpha_2 \xi_2 + \alpha_3 \xi_3 \equiv 0 \ (\mathrm{mod}\ 2)}
%CONT>>
